% !TeX root = ../main.tex
\section{Introduction}\label{sec:intro}

Single-instruction-multiple-data (SIMD) units let one instruction operate on a whole vector of
elements. On the Intel Core i9-10900F studied here the widest available unit is AVX2 with FMA, a
\qty{256}{bit} register that holds \num{8} float32 or \num{4} float64 lanes, so the \emph{ideal}
speedup over scalar code is $8\times$ (float32) or $4\times$ (float64). Real kernels rarely reach
that ceiling: loads and stores, loop prologue/epilogue handling, reduction dependencies, and---above
all---finite memory bandwidth cap the benefit. The purpose of this report is to measure where the
benefit appears, by how much it falls short of the lane count, and exactly when it disappears.

This report answers:
\begin{enumerate}
  \item How much faster are auto-vectorized builds than scalar builds for each kernel, and how does
        that change from cache-resident to DRAM-resident working sets? (\cref{sec:baseline,sec:locality})
  \item What do alignment, loop tails, stride and gather access, and data type cost? (\cref{sec:alignment,sec:stride,sec:dtype})
  \item Where do these kernels sit in a roofline model of this machine? (\cref{sec:roofline})
  \item When does the expected SIMD benefit fail to appear? (\cref{sec:negative})
\end{enumerate}

In brief: SIMD helps most when data fits in L1/L2, where \code{native} gives roughly $5.8\times$ for
SAXPY and the cache feeds the vector units fast enough to approach the lane count; the advantage
shrinks step-by-step as the working set crosses L2 and L3 and falls to $1.88\text{--}1.99\times$ for SAXPY
in DRAM ($1.5\text{--}1.6\times$ multiply, $2.6\text{--}2.8\times$ stencil), where the vector build runs at the memory-bandwidth roof and only the scalar build's own
instruction cost keeps the ratio above $1\times$. Data type matters through lane count (float64
roughly halves the gain) and access pattern matters even more (any stride~$\ge2$ removes the SIMD
advantage; gather, even with AVX2 \code{vgatherdps}, never gains). Two cases show the limits of ``it vectorized, so it
must be faster'': memory-bound SAXPY (\cref{sec:negative}) and the dot-product reduction
(\cref{sec:vectorization}).
